\documentclass[UTF8,twocolumn,10pt,a4paper]{ctexart}

\usepackage[a4paper,left=1.55cm,right=1.55cm,top=1.45cm,bottom=1.60cm,columnsep=0.66cm]{geometry}
\usepackage{fontspec}
\setmainfont{DejaVu Serif}
\setsansfont{DejaVu Sans}
\setmonofont{DejaVu Sans Mono}
\IfFontExistsTF{Noto Serif CJK SC}{\setCJKmainfont{Noto Serif CJK SC}}{}
\IfFontExistsTF{Noto Sans CJK SC}{\setCJKsansfont{Noto Sans CJK SC}}{}
\IfFontExistsTF{Noto Sans Mono CJK SC}{\setCJKmonofont{Noto Sans Mono CJK SC}}{}

\usepackage{amsmath,amssymb,mathtools,bm}
\usepackage{booktabs,tabularx,array}
\usepackage{enumitem}
\usepackage{microtype}
\usepackage{xcolor}
\usepackage[hidelinks]{hyperref}

\hypersetup{
  pdftitle={BRIAN RC-KV: Implemented Reader-Compiled KV Cache and CPBC Prefill},
  pdfauthor={YMH-Latent-Sphere},
  pdfsubject={Implementation-grounded report for reader-compiled KV cache under free routing},
  pdfkeywords={BRIAN, free routing, RC-KV, BDRE, CPBC, KV cache, TBPTT}
}

\setlength{\parindent}{1em}
\setlength{\parskip}{0.12em}
\renewcommand{\baselinestretch}{1.015}
\setlength{\abovedisplayskip}{4.4pt plus 1pt minus 1pt}
\setlength{\belowdisplayskip}{4.4pt plus 1pt minus 1pt}
\setlength{\abovedisplayshortskip}{2.4pt plus 1pt}
\setlength{\belowdisplayshortskip}{2.4pt plus 1pt minus 1pt}
\setlength{\textfloatsep}{7pt plus 2pt minus 2pt}
\setlength{\dbltextfloatsep}{7pt plus 2pt minus 2pt}
\setlength{\intextsep}{6pt plus 2pt minus 2pt}
\setlength{\emergencystretch}{1.5em}
\ctexset{
  section={beforeskip=0.72ex plus .2ex,afterskip=0.34ex plus .1ex,format=\large\bfseries},
  subsection={beforeskip=0.58ex plus .2ex,afterskip=0.26ex plus .1ex,format=\normalsize\bfseries},
  paragraph={beforeskip=0.45ex plus .15ex,afterskip=0.55em,runin=true,format=\normalsize\bfseries}
}
\setlist{nosep,leftmargin=1.45em}
\allowdisplaybreaks

\newcolumntype{Y}{>{\raggedright\arraybackslash}X}
\newcommand{\Concat}{\operatorname{Concat}}
\newcommand{\TopK}{\operatorname{TopK}}
\newcommand{\softmax}{\operatorname{softmax}}
\newcommand{\RCKV}{\textsf{RC-KV}}
\newcommand{\BDRE}{\textsf{BDRE}}
\newcommand{\CPBC}{\textsf{CPBC}}
\newcommand{\DP}{\textsf{CPBC-DP}}
\newcommand{\FB}{\textsf{CPBC-FB}}
\newcommand{\Serial}{\textsf{BDRE-Serial}}
\newcommand{\code}[1]{\texttt{#1}}

\title{\vspace{-1.0em}\textbf{BRIAN--\RCKV{}：自由路由下的 Reader-Compiled KV 缓存\\与 Chunkwise Progressive Bank Completion}}
\author{YMH-Latent-Sphere}
\date{实现报告，2026 年 7 月 16 日}

\begin{document}
\maketitle
\vspace{-1.55em}

\begin{abstract}
自由路由使不同 token 在同一计算深度访问不同 block，并允许重复访问同一 block；因此，固定层 Transformer 中以 layer identity 为索引的 KV cache 无法直接迁移到 BRIAN。本文给出已经实现的 Reader-Compiled KV Cache（\RCKV{}）系统，而非早期理论规格的复述。系统首先由每个 writer block 独立地将其多头 K/V 压缩到 canonical latent，再由 Block--Depth Relative Encoding（\BDRE{}）按照 reader block、可选 reader step 及球面几何编译 reader-specific cache；reader head 使用独立映射恢复 Key 视角，并在压缩 Value 空间中完成等价聚合。实现同时提供精确 token-serial oracle、三种 current-token self-KV、三类 reader-step cache policy、可选 hard top-$k$ 与独立 K/V 温度。

为解决精确 teacher-forced 训练的串行瓶颈，本文进一步定义 Chunkwise Progressive Bank Completion（\CPBC{}）：chunk 内 token 并行、route depth 串行，cache bank 随 recurrent step 渐进补全，并以严格 token-causal mask 阻止未来信息。实现包含 depth-prefix visibility（\DP{}）和 full-bank visibility（\FB{}）、跨 chunk 的 stateful TBPTT，以及 rank-local state 下的显式 DDP 梯度同步。完整测试验证了因果性、incremental consistency、cache policy 等价性和分布式更新一致性。grouped expert GEMM、FlexAttention fused reader、bounded reader batching 与向量化 FB compiler 将两卡 \DP{} C2048 smoke 提升到约 $4.55\times10^4$ token/s；匹配的普通 Transformer 约为 $6.19\times10^5$ token/s。该优化没有创建新模型版本，剩余差距主要来自动态分组、route-dependent padding、attention 算术与高频 dispatch。
\end{abstract}

\noindent\textbf{关键词：} BRIAN；自由路由；Reader-Compiled KV；BDRE；CPBC；TBPTT

\section{引言}
固定层 decoder-only Transformer~\cite{vaswani2017attention} 将第 $\ell$ 层写出的历史 Key/Value 交给未来 token 的第 $\ell$ 层读取。该机制隐含一个稳定假设：writer identity 与 reader identity 都由层号 $\ell$ 决定。BRIAN 的 free routing 打破了此假设。token 可以选择不同的 free-block 序列、在路径中跳过 block、重复访问 block，并在不同步数退出。此时把 block $b$ 写出的 K/V 原样暴露给未来任意 reader，不仅混合了不同 head 的坐标系，也把历史路径长度扩展为额外的 attention object 轴。

\RCKV{} 的目标不是压缩计算量，也不是迫使路径变短，而是建立一个适用于自由路由的记忆接口：writer 保留自己的参数化，缓存进入跨路径可组合的 canonical space，reader 仍保留 block- 和 head-specific 的解释能力。本文以当前代码、配置、测试和 B200 实测为准，并对早期理论稿作出以下实现性修正：

\begin{enumerate}
  \item 默认 persistent cache 仅按 reader block 编译；reader-step 维度是可选语义，而非强制的 depth--block 笛卡尔积。
  \item free blocks 的 $W_Q,W_K,W_V$、writer canonicalizer 和 reader decoder 均不共享；``shared'' 只表示 canonical cache interface 共享。
  \item temporal RoPE 在 decoded Key 上显式执行；只有 Value 路径使用严格等价的低维聚合。
  \item 精确 token-by-token 执行永久保留为因果性与 incremental inference oracle。
  \item 并行训练采用独立的 \CPBC{} 语义，并显式报告 visibility policy 与梯度截断边界。
\end{enumerate}

本文不声称当前正式 5B 训练已经证明下游能力提升。本文的贡献是给出真实系统合约、已验证性质、可选机制和剩余性能瓶颈。

\section{模型范围与符号}
\subsection{当前 R125 实例}
当前正式实例含两个固定 pre blocks、八个 free route blocks、一个 Exit Block 和两个固定 post blocks。固定 pre/post blocks 具有稳定 layer identity，继续使用标准逐层 autoregressive KV cache；\RCKV{} 只作用于八个 free blocks。表~\ref{tab:scope} 给出已解析配置和实际模型统计。

\begin{table}[t]
\centering
\caption{当前 BRIAN-R125--\RCKV{} 实例。}
\label{tab:scope}
\small
\begin{tabular}{lr}
\toprule
项目 & 数值 \\
\midrule
参数量 & $140{,}308{,}105$ \\
hidden dimension $d$ & 768 \\
attention heads $H$ & 12 \\
head dimension $d_h$ & 64 \\
fixed pre/free/fixed post & $2/8/2$ \\
最大 route steps $S_{\max}$ & 16 \\
route action & hard top-1 \\
position dimension $p$ & 64 \\
canonical Key/Value & $r_K=r_V=32$ \\
context length & 2048 \\
\bottomrule
\end{tabular}
\end{table}

令 token $i$ 的有效 free route 为
\begin{equation}
  \pi_i=(b_{i,0},b_{i,1},\ldots,b_{i,S_i-1}),
  \qquad S_i\le S_{\max},
  \label{eq:route}
\end{equation}
其中 $b_{i,s}\in\{0,\ldots,R-1\}$，当前实现 $R=8$。OUT 是第 $R$ 个 route action，但不产生 writer code。用 $v_{i,s}\in\{0,1\}$ 表示 step $s$ 是否为有效 internal write；已退出 token 的 hidden state 不再更新，后续位置仅作为 padding/OUT 处理。

\subsection{球面位置几何}
系统维护 $\bm z_{\mathrm{IN}}$、八个 $\bm z_b$ 和 $\bm z_{\mathrm{OUT}}$。初始化时 IN/OUT 互为对跖点，internal blocks 在与该轴正交的子空间中形成 regular simplex：
\begin{align}
  \bm z_{\mathrm{IN}}^\top\bm z_{\mathrm{OUT}}&=-1,\\
  \bm z_{\mathrm{IN}}^\top\bm z_b
  =\bm z_{\mathrm{OUT}}^\top\bm z_b&=0,\\
  \bm z_b^\top\bm z_{b'}&=-\frac{1}{R-1},\quad b\ne b'.
  \label{eq:spherical-code}
\end{align}
构造后施加固定随机正交旋转。所有 position 可学习，并在读取时归一化。IN 初始化 router 与第一步 position adapter；internal positions 同时服务于 router、block adapter 和 \BDRE{}；OUT 只表示终止位置，不索引 reader cache。权重 $10^{-3}$ 的 Gram-matrix loss 弱保持式~\eqref{eq:spherical-code}，但不构造邻近路径或模仿路径。

\section{Reader-Compiled KV Cache}
\subsection{Block-specific writer canonicalization}
token $i$ 在 step $s$ 进入 block $b=b_{i,s}$。令 $\bm h_{i,s}$ 为路由前 hidden state，$g_b(\cdot)$ 表示现有 position adapter，则标准 pre-norm attention 输入为
\begin{equation}
  \bm x_{i,s}=\operatorname{RMSNorm}
  \bigl(\bm h_{i,s}+g_b(\bm z_{i,s})\bigr).
\end{equation}
每个 block 的 packed QKV linear 在数学上仍对应独立的 block/head 参数：
\begin{align}
  \bm q_{i,s}^{(h)}&=W_{Q,b}^{(h)}\bm x_{i,s},\\
  \bm k_{i,s}^{(h)}&=W_{K,b}^{(h)}\bm x_{i,s},\\
  \bm v_{i,s}^{(h)}&=W_{V,b}^{(h)}\bm x_{i,s}.
\end{align}
在施加 temporal RoPE 之前，全部 Key heads 和 Value heads 分别沿 feature 轴拼接，并由当前 writer block 的独立矩阵压缩：
\begin{align}
  \bm c^K_{i,s}
  &=A^K_b\Concat_{h=1}^{H}\bm k_{i,s}^{(h)}
  \in\mathbb R^{r_K},\\
  \bm c^V_{i,s}
  &=A^V_b\Concat_{h=1}^{H}\bm v_{i,s}^{(h)}
  \in\mathbb R^{r_V}.
  \label{eq:writer-code}
\end{align}
实现中的 $A^K_b,A^V_b$ 是每个 free block 独立的无 bias linear。它们不在 block 间共享；canonical space 的可组合性由训练学习，而不是由参数 tying 强制获得。

\subsection{BDRE compiler}
对 reader block $r$、可选 reader step $\ell$ 与 writer step $s$，实现使用统一分数
\begin{equation}
  e_{i,\ell,r,s}
  =\tau_z\,\bar{\bm z}_r^\top\bar{\bm z}_{b_{i,s}}
  -\lambda_s d(\ell,s)
  +\gamma\frac{s+1}{S_{\max}},
  \label{eq:bdre-score}
\end{equation}
其中 $\bar{\bm z}=\bm z/\lVert\bm z\rVert_2$。$d(\ell,s)$ 可取 $|\ell-s|$，当前 \CPBC{} profile 使用归一化形式
\begin{equation}
  d(\ell,s)=\frac{|\ell-s|}{S_{\max}-1}.
\end{equation}
三项均可独立退化：block-only mode 令 $\lambda_s=0$ 且不需要 $\ell$；late-step bias 由 $\gamma>0$ 开启，默认 $\gamma=0$；position contribution 由 $\tau_z$ 控制。

设 $\mathcal V_{i,\ell}$ 是当前语义允许读取的有效 writer steps。hard compile top-$k_c$ 是可选项：
\begin{equation}
  \mathcal I_{i,\ell,r}=
  \begin{cases}
    \mathcal V_{i,\ell}, & k_c=\varnothing,\\
    \TopK_{s\in\mathcal V_{i,\ell}}(e_{i,\ell,r,s};k_c),
      & \text{otherwise}.
  \end{cases}
\end{equation}
Key 与 Value 共享 support，但采用独立温度：
\begin{align}
  \alpha^X_{i,\ell,r,s}
  &=\frac{\exp(e_{i,\ell,r,s}/T_X)}
  {\sum_{u\in\mathcal I_{i,\ell,r}}
    \exp(e_{i,\ell,r,u}/T_X)},\\
  \bm m^X_{i,\ell,r}
  &=\sum_{s\in\mathcal I_{i,\ell,r}}
    \alpha^X_{i,\ell,r,s}\bm c^X_{i,s},
  \qquad X\in\{K,V\}.
  \label{eq:compiled-cache}
\end{align}
softmax 在 FP32 中执行后再转换回模型 dtype。当前默认 $T_K=0.5$、$T_V=1.0$、$k_c=\varnothing$，因而 Key fusion 比 Value fusion 更尖锐，但所有有效 writer steps 都保留梯度。

\subsection{Persistent state variants}
实际实现支持表~\ref{tab:cache-modes} 中三类历史状态。block-only 是精确 reference 的最小 persistent form；reader-step 模式可 eager 编译、首次读取时 lazy 编译并 memoize，或保留 writer history 动态重编译。\CPBC{} 强制使用 eager step cache。

\begin{table*}[t]
\centering
\caption{已经实现的 persistent cache 组织。$B,T,L,R$ 分别为 batch、token、reader-step 和 free-block 维。}
\label{tab:cache-modes}
\small
\begin{tabularx}{\textwidth}{l l l Y}
\toprule
模式 & K/V 主要形状 & R125 BF16 每 token & 实现语义 \\
\midrule
block-only & $[B,T,R,r_X]$ & K+V 共 1 KiB & 每个 completed token、每个 reader block 一对 cache \\
reader-step eager & $[B,T,L,R,r_X]$ & K+V 共 16 KiB & token 完成时编译全部 reader steps \\
reader-step lazy & block cache + writer history + memo & 按访问增长 & 首次访问 $(r,\ell)$ 时编译整段历史并缓存 \\
reader-step dynamic & $[B,T,S_{\max},r_X]$ writer history & 约 2 KiB & 每次读取重新融合，仅用于 correctness/cost reference \\
\bottomrule
\end{tabularx}
\end{table*}

在 $T=2048$ 时，reader-step eager 的 canonical K/V 约为每个 sequence 32 MiB；该数字不含 pre/post KV、activation、mask、梯度与 allocator overhead。当前 formal \CPBC{} 使用这一 step-indexed state。

\section{Reader Attention 与 Current-Token Memory}
\subsection{Head-specific read}
每个 reader block $r$ 和 head $h$ 拥有独立参数
\begin{align}
  B^K_{r,h}&\in\mathbb R^{r_K\times d_h},\\
  B^V_{r,h}&\in\mathbb R^{r_V\times d_h}.
\end{align}
代码中它们分别存为形状 $[H,r_K,d_h]$ 与 $[H,r_V,d_h]$ 的参数。对历史 token $i$，显式 decoded Key 为
\begin{equation}
  \widehat{\bm k}_{i,\ell,r}^{(h)}
  =(B^K_{r,h})^\top\bm m^K_{i,\ell,r}.
\end{equation}
query 和 decoded Key 按真实 token position 施加 RoPE~\cite{su2021roformer}：
\begin{align}
  \widetilde{\bm q}_{t,\ell,r}^{(h)}
  &=\operatorname{RoPE}(\bm q_{t,\ell,r}^{(h)},t),\\
  \widetilde{\bm k}_{i,\ell,r}^{(h)}
  &=\operatorname{RoPE}(\widehat{\bm k}_{i,\ell,r}^{(h)},i),\\
  a_{t,i}^{(\ell,r,h)}
  &=\softmax_{i\le t}
  \left(\frac{(\widetilde{\bm q}_{t,\ell,r}^{(h)})^\top
  \widetilde{\bm k}_{i,\ell,r}^{(h)}}{\sqrt{d_h}}\right).
  \label{eq:reader-attn}
\end{align}

Value 路径使用严格等价的低维聚合：
\begin{align}
  \bm u_{t,\ell,r}^{(h)}
  &=\sum_{i\le t}a_{t,i}^{(\ell,r,h)}
  \bm m^V_{i,\ell,r},\\
  \bm o_{t,\ell,r}^{(h)}
  &=(B^V_{r,h})^\top\bm u_{t,\ell,r}^{(h)}.
  \label{eq:latent-value}
\end{align}
式~\eqref{eq:latent-value} 只把线性 decode 移到求和之后，不改变注意力定义。Key 路径不能直接使用 $(B^K)\bm q$ 与 latent Key 点积，因为每个历史 Key 具有不同 RoPE rotation；当前唯一允许的 Key read mode 是 \code{explicit\_decode\_rope}。

\subsection{当前 token 的三种 self-KV}
精确 token-serial backend 支持：
\begin{itemize}
  \item \code{bdre\_prefix}：把当前 step code 追加到本 token 的 route prefix，并按当前 reader 编译为一个 self object；
  \item \code{current\_step}：只使用当前 step 的 $\bm c^K,\bm c^V$；
  \item \code{none}：只允许 $i<t$ 的历史 token。
\end{itemize}
默认 \code{bdre\_prefix} 不产生循环依赖，因为当前 K/V 由 attention input 在 attention output 之前计算。它也不会把同一 token 的每个 writer step 作为独立 temporal object 暴露给 attention。

在 \CPBC{} 中，一个 active chunk 的每个 token 在 reader step $\ell$ 都产生一个 on-the-fly prefix-compiled object；严格的 token mask 只允许 query $t$ 读取 $i\le t$。因此 current-token self object 与同 chunk 先前 token 使用同一 bank 表示，不会重复展开 writer-step 轴。

\section{精确执行：\Serial{}}
\Serial{} 是 sequence-serial、batch-parallel 的语义 oracle。对 token position $t$，实际执行顺序为：
\begin{enumerate}
  \item fixed pre blocks 使用各自的标准 incremental KV 处理 token $t$；
  \item route position 初始化为 IN，router 在每个 step 为 batch 内 token 选择一个 internal block 或 OUT；
  \item active token 按选中 block 分组，读取 $i<t$ 的 reader-compiled cache，并按 self-KV mode 加入当前 token；
  \item 每个有效 step 记录式~\eqref{eq:writer-code} 的 writer code；退出 token 的 hidden state 停止更新；
  \item 完整 route 结束后才把 token $t$ 编译并提交到 persistent cache；
  \item Exit Block 与 fixed post blocks 处理 token $t$，随后产生 LM logits。
\end{enumerate}

完整 teacher-forced \code{forward} 逐 token 调用同一个 incremental kernel；\code{forward\_incremental} 则显式接收并返回 pre/post KV 与 \RCKV{} state，不重算 prefix。测试中两者的 logits 最大差为 0。由于 token $t$ 的 persistent cache 只在 route 完成后提交，未来 suffix 的替换无法改变 prefix logits、router logits 或 cache state。

精确 reader-step mode 对 completed token 可采用 eager/lazy/dynamic 三种物理策略；它们的数学输入均是该 token 的完整有效 route。该语义与后文 chunk 内渐进补全并不总是相同。

\section{CPBC：近似 Prefill}
\subsection{定义与执行}
精确 backend 在一张 B200 上仅约 25--28 token/s，无法用于 5B-token 训练。\CPBC{} 将长度 $T$ 的序列切为连续 chunks $\mathcal C_0,\ldots,\mathcal C_{M-1}$，每个 chunk 最多含 $C$ 个 token。它在 token 轴并行，在 route-depth 轴保持串行。

对 chunk $\mathcal C_m$ 与 reader step $\ell$：
\begin{enumerate}
  \item fixed pre blocks 对整段 chunk 执行标准 causal attention；
  \item router 独立地为 $[B,C]$ 中每个 active token 选择 action；
  \item token 按 reader block 分组计算 Q/K/V 与 canonical writes；
  \item compiler 只用当前已生成的 $s\le\ell$ writer prefix，为所有 reader blocks 生成 step-$\ell$ bank；
  \item selected reader 同时读取 completed history 与当前 chunk bank，并施加 $i\le t$ 的绝对 token-causal mask；
  \item 已退出 token 冻结，其余 token 进入 $\ell+1$；最后持久化 bank，经 fixed post blocks 输出 logits。
\end{enumerate}

当前 chunk 的合法 writer support 为
\begin{equation}
  \mathcal V^{\mathrm{active}}_{i,\ell}
  =\{s\mid v_{i,s}=1,\ s\le\ell\}.
  \label{eq:active-prefix}
\end{equation}
式~\eqref{eq:active-prefix} 是 hard mask，发生在 compiler softmax 之前；修改 $s>\ell$ 的 writer code 不会改变 earlier reader-step cache。

\subsection{Depth Visibility Policy}
\CPBC{} 的近似来自 completed chunks 对 route-depth bank 的可见性定义。实现提供两种 policy。

\paragraph{Depth-Prefix Visibility（\DP{}）。}
active 与 completed token 都保持
\begin{equation}
  \mathcal V^{\mathrm{DP}}_{i,\ell}
  =\{s\mid v_{i,s}=1,\ s\le\ell\}.
  \label{eq:dp}
\end{equation}
因此在固定权重下，改变 chunk boundary 不改变 \DP{} forward value；chunk size 仍会改变 activation shape 与梯度边界。

\paragraph{Full-Bank Visibility（\FB{}）。}
active chunk 仍使用式~\eqref{eq:active-prefix}，但 chunk 完成后为每个 reader step $\ell$ 使用全部有效 writers 重编译：
\begin{equation}
  \mathcal V^{\mathrm{FB}}_{i,\ell}
  =\{s\mid v_{i,s}=1\},\qquad i\in\mathcal C_{<m}.
  \label{eq:fb}
\end{equation}
reader-step $\ell$ 仍进入式~\eqref{eq:bdre-score}，``full bank'' 只表示 candidate support 包含 completed route。\FB{} 的 forward 依赖 chunk boundary；当 $C=1$ 时，它在已测试配置中与 token-serial completed-route semantics 匹配到 $2\times10^{-5}$ 容差。

\begin{figure*}[t]
\centering
\fbox{\parbox{0.96\textwidth}{\small
\textbf{CPBC bank lifecycle.}
对 active chunk，route step $\ell$ 只产生 prefix bank $s\le\ell$；同 chunk attention 通过 token-causal mask 读取这些 bank。chunk 完成后，\DP{} 原样持久化每个 depth prefix，\FB{} 则用完整有效 route 为所有 reader steps 重编译。未来 chunk 只读取已持久化 bank。TBPTT detach 改变 autograd history，不改变这些 forward cache values。}}
\caption{\CPBC{} 中 cache bank 的渐进补全、promotion 与梯度边界。}
\label{fig:cpbc-life}
\end{figure*}

\section{Stateful TBPTT 与 DDP}
\subsection{梯度跨度}
令训练 chunk size 为 $C$，每经过 $U$ 个 chunks 才 detach persistent state。最大梯度跨度为
\begin{equation}
  H_{\mathrm{grad}}=\min(T,C U).
  \label{eq:grad-horizon}
\end{equation}
同一组内的 chunk losses 相加后执行一次 backward；随后 pre/post KV 与 \RCKV{} bank 保留数值但切断 autograd graph。所有 chunks 完成后才执行 optimizer step，避免一个 sequence 同时依赖两个参数版本。

LM loss 以完整 sequence 的有效 next-token label 数归一化：
\begin{equation}
  \mathcal L_{\mathrm{LM}}
  =\frac{1}{B(T-1)}\sum_{n=1}^{B}\sum_{t=0}^{T-2}
  -\log p_\theta(x_{n,t+1}\mid x_{n,\le t}).
\end{equation}
因此 chunk losses 的和与完整 sequence mean loss 一致。$U$ 在固定权重下不改变 forward value，但决定 later chunk loss 能否更新 earlier cache writers。

\subsection{Rank-local state 与手工同步}
DDP 仅划分 batch，每个 rank 维护自己的 cache state。所有 stateful forwards 在 \code{no\_sync()} 下运行，完整 optimizer-step accumulation 后进行一次显式同步：先 all-reduce parameter-used mask，再把全局使用但本 rank 未使用的参数补零梯度，最后按 dtype/device bucket all-reduce 并除以 world size。当前两卡 formal profile 每步同步约 561 MB gradients，分为 9 个 64 MB 上限 buckets。CPU merged-batch 测试验证了 loss、gradient 与 parameter update 等价。这里的跨 chunk 截断属于 TBPTT~\cite{williams1990tbptt}，但 cache state 与稀疏参数同步是本实现的显式扩展。

\subsection{当前正式训练 profile}
表~\ref{tab:formal-profile} 是正在运行的实现 profile，而不是所有 \RCKV{} mode 的默认值。

\begin{table}[t]
\centering
\caption{当前 formal \DP{} profile。}
\label{tab:formal-profile}
\small
\begin{tabular}{ll}
\toprule
设置 & 数值 \\
\midrule
execution & \code{synchronous\_prefix} \\
visibility & \code{depth\_prefix} \\
attention backend & \code{shared\_padded\_explicit} \\
dispatch & \code{grouped\_mm} \\
$C,U$ & $512,1$ \\
DDP & 2 B200 \\
local/global batch & $16/32$ \\
tokens/update & 65,536 \\
precision & BF16 \\
learning rate & $3\times10^{-4}$ \\
training budget & 5B tokens \\
\bottomrule
\end{tabular}
\end{table}

当前 routing curriculum 在 step 0--1500 将 router probability 从 0 逐步升至 1，但 route-imitation loss 始终为 0。logit noise 从 0.08 缓慢衰减至 0.008，random-route override 从 0.25 衰减至 0.03。约束包括 min exit step 4、step 12 开始的 exit ramp、最终强制 OUT 和 self-recurrence cap 2。当前 objective 为
\begin{align}
  \mathcal L={}&\mathcal L_{\mathrm{LM}}
  +2\!\times\!10^{-4}\mathcal L_{\mathrm{cost}}
  +0.02\mathcal L_{\mathrm{selected}}
  +0.05\mathcal L_{\mathrm{coverage}}\nonumber\\
  &+0.02\mathcal L_{\mathrm{exit}}
  +0.001\mathcal L_{\mathrm{geometry}}.
  \label{eq:objective}
\end{align}
location、plain balance、transition-diversity 与 route-imitation 权重为 0。该 routing recipe 是当前实验控制，不是 \RCKV{} 数学定义的一部分。

在不改变上述 formal run 的前提下，当前还准备了同一 \RCKV{} 定义下的 fused-reader profile。给定相同 route decisions，\DP{} cache/attention 满足 chunk-boundary invariance。C2048 将 U1 gradient horizon 从 512 提高到 2,048 token，并把四次 backward group 合为一次；scheduled routing、noise 与 random override 还会按不同 tensor grouping 消耗 RNG，因此训练轨迹不与 C512 等价。

\section{已实现选项与合法组合}
表~\ref{tab:options} 总结当前代码可以选择的机制。部分选项是 correctness reference 或 ablation，不应视为推荐生产后端。

\begin{table*}[t]
\centering
\caption{\RCKV{}/\CPBC{} 的真实配置表。``Formal'' 指表~\ref{tab:formal-profile} 的当前运行。}
\label{tab:options}
\scriptsize
\setlength{\tabcolsep}{3pt}
\begin{tabularx}{\textwidth}{p{0.15\textwidth} p{0.25\textwidth} p{0.14\textwidth} Y}
\toprule
配置 & 已实现取值 & Formal & 语义/约束 \\
\midrule
Execution mode & token-by-token；synchronous-prefix & synchronous-prefix & 前者为 oracle，后者为 CPBC \\
Depth mode & none；reader-step；synchronous-prefix & synchronous-prefix & CPBC 必须与 synchronous-prefix depth mode 配对 \\
Visibility policy & depth-prefix；full-bank & depth-prefix & full-bank 仅允许 CPBC \\
Reader cache & eager；lazy；dynamic & eager & CPBC 强制 eager；其余服务 exact reader-step mode \\
Self-KV & bdre-prefix；current-step；none & bdre-prefix & CPBC 当前强制 bdre-prefix \\
Attention backend & per-query-reference；shared-padded-explicit；shared-padded-flex & shared-padded-explicit & Flex 路径融合 decode、RoPE、causal attention 与 Value read \\
Flex reader group & $1\ldots B$ & N/A & 只改变 reader 执行布局；C512/FB-C128 用 8，C2048 用 1 \\
FB compiler & loop；vectorized & loop & vectorized 同时计算所有 completed reader depths \\
Dispatch & legacy-cuda-scan；grouped-host；grouped-mm & grouped-mm & grouped-mm 批处理 routed expert projections \\
Compiler top-$k$ & null 或正整数 & null & Key/Value 共用 support；不得超过 $S_{\max}$ \\
$T_K,T_V$ & 正实数 & $0.5,1.0$ & 独立 compiler softmax 温度 \\
$\tau_z,\lambda_s,\gamma$ & 非负实数 & $1.0,0.25,0$ & formal depth distance 归一化 \\
Detach interval & $U\ge1$ & 1 & 只改变梯度跨度，不改变固定权重 forward \\
\bottomrule
\end{tabularx}
\end{table*}

实现还设置了硬 guardrails：\RCKV{} v1 只允许 route top-1；\CPBC{} 要求 hard exit；position mode 必须为 spherical code 且 IN position 独立；不得与旧 hidden Global KV、attention-summary Global KV 或 parallel passing 同时启用。Top-2 与 weighted route fusion 尚未纳入该版本。

\section{正确性验证}
表~\ref{tab:correctness} 汇总已经执行的核心验收。所有比较使用同参数、同输入和同随机条件；BF16 比较使用显式容差。

\begin{table*}[t]
\centering
\caption{当前实现的 correctness evidence。}
\label{tab:correctness}
\small
\begin{tabularx}{\textwidth}{Y l}
\toprule
检查 & 结果 \\
\midrule
Exact full forward 与 stateful incremental & logits 最大差 $0$ \\
Exact streamed chunks 与未分块 exact forward & forward/loss 一致 \\
reader-step eager/lazy/dynamic & $\mathrm{atol}=\mathrm{rtol}=10^{-6}$ 内一致 \\
prefix suffix invariance 与 future-writer intervention & 通过；future token/step 不影响 prefix \\
\FB{}，chunk 1 与 \Serial{} & reference/shared-explicit 均在 $2\times10^{-5}$ 内一致 \\
\DP{} cache/attention chunk-boundary invariance & 相同 route decisions 下通过 \\
CPBC full forward、stream chunks 与 one-token incremental & 在 $2\times10^{-5}$ 内一致 \\
$U=1$ 与 $U=2$ & forward loss 相同，gradient 按预期不同 \\
shared-padded-explicit 与 per-query-reference & logits、loss、cache 与 parameter gradients 一致 \\
vectorized FB compiler 与 reader-depth loop & output、weights 与 writer/position gradients 一致 \\
CUDA Flex DP/FB suffix 与 streamed forward & 通过 \\
explicit 与 Flex BF16 reader & reported loss 相同；logits/gradients 在 BF16 舍入范围 \\
CPU DDP 与 merged global batch & loss、gradient、update 一致 \\
rank-local unused parameter synchronization & 通过 \\
B200 CUDA BF16 forward/backward 与 DDP2 smoke & finite 且 checkpoint state 通过 \\
完整 repository regression & 606 passed，15 个既有 warnings \\
\bottomrule
\end{tabularx}
\end{table*}

上述结果证明实现合约与因果接口成立，但不证明 \DP{} 或 \FB{} 在语言建模与 public benchmark 上优于 baseline。模型选择必须同时观察 PPL、reasoning/public suite、route entropy、path diversity、block coverage、compiler entropy、writer mass 和 cache norm。

\section{性能结果与瓶颈}
\subsection{Kernel/backend 演进}
表~\ref{tab:performance} 给出在 B200 上完成的可复现实测。早期 token-serial 只适合作为 oracle；\code{shared\_padded\_explicit} 把同一 $(\mathrm{reader},\ell,\mathrm{batch})$ 的 queries 打包，使 decoded history 在组内共享，并在 latent Value 上聚合。它相对 per-query reference 提升约 $3.2\times$，同时显著降低 activation memory。

\begin{table*}[t]
\centering
\caption{R125 BF16 B200 性能。标记 DDP2 之外均为单卡完整 forward+backward、包含 warmup 后的测量。}
\label{tab:performance}
\small
\begin{tabular}{lrrrrrr}
\toprule
Backend & BS & $T$ & $C$ & token/s & peak allocated & 备注 \\
\midrule
Exact token-serial & 32 & 128 & 8 & 28.18 & 3.28 GiB & oracle \\
Exact token-serial & 32 & 256 & 8 & 25.53 & 4.46 GiB & oracle \\
CPBC per-query reference & 8 & 2048 & 256 & 1,467.80 & 111,353 MiB & median \\
CPBC shared-padded-explicit & 8 & 2048 & 256 & 4,695.64 & 27,333 MiB & median \\
CPBC shared-padded-explicit & 14 & 2048 & 512 & 11,658.95 & 101,808 MiB & stable candidate \\
CPBC grouped-MM & 16 & 2048 & 512 & 15,167 & 120,947 MiB & accepted single-B200 shape \\
Formal \DP{} grouped-MM DDP2 & 32 global & 2048 & 512 & 28,638--30,609 & 124.5 GiB/rank & accepted smoke \\
\DP{} fused reader, group 8 & 16 & 2048 & 512 & 19,139 & 24,742 MiB & same model definition \\
\DP{} fused reader, group 1 & 16 & 2048 & 2048 & 23,246 & 66,011 MiB & horizon 2048 \\
\FB{} fused/vectorized, group 8 & 16 & 2048 & 128 & 9,013 & 13,460 MiB & $>10\times$ historical FB \\
\DP{} C2048 fused DDP2 & 32 global & 2048 & 2048 & 45,443--45,648 & 70.2 GiB/rank & steady smoke \\
\bottomrule
\end{tabular}
\end{table*}

两个 DDP2 profile 均保持 global batch 32。accepted C512 grouped-MM 的两个 steady steps 平均约 29,624 token/s；prepared C2048 fused profile 平均约 45,546 token/s，约为单卡的 $1.94\times$。同期匹配数据与 global batch 的普通 Transformer 平均为 618,738 token/s，因此系统级差距仍约 $13.6\times$。C2048 相比最终 fused C512 的单卡提升约 21.5\%，而不是 chunk 变为四倍所暗示的 $4\times$；attention 工作量和 padding 阻止线性扩展。

\subsection{剩余热点}
profile 与实现检查显示主要瓶颈为：
\begin{enumerate}
  \item 每个 route step 把 selected actions 从 GPU 同步到 CPU，并由 Python 构造 active block groups；
  \item grouped GEMM 已合并主要 expert projections，但 reader packing、copy 与 index kernels 仍高度碎片化；
  \item C2048 下 fused attention forward/backward 已成为主要 GPU 算术；
  \item reader action 间 token 数不均，跨 reader batching 会引入 padding；实测 C2048 的 group 2/8 均慢于 group 1；
  \item $C=512,U=1,T=2048$ 每个 optimizer step 执行四次 backward groups，随后同步约 561 MB gradients；
  \item step-indexed bank 与 route-dependent padding 增加 memory traffic，且不同路径长度使 group shape 高频变化。
\end{enumerate}
compiler 本身在当前运行中通常只占少量 wall time；单纯删减 compiler score 或缩短 route 不能解决主要工程问题，也会改变研究目标。grouped expert GEMM、窄 cache gather、fused reader、bounded reader batching 与 vectorized FB compiler 已实现；下一阶段性能工作应优先验证 GPU-resident packing 和 routed RMSNorm fusion，并以 logits/loss/gradient 等价为验收条件。

\section{局限与尚未回答的问题}
\paragraph{训练速度。}
当前 \CPBC{} 已把不可训练的 token-serial 速度提高数百倍，但仍比 matched baseline 慢约一个数量级以上。完整 5B anchor 的墙钟成本仍高，阻碍大规模 ablation。

\paragraph{CPBC 是显式近似语义。}
\DP{} 对 chunk boundary 前向不变，但它让 completed token 的 reader depth $\ell$ 永久只看 writer prefix $s\le\ell$；\FB{} 更接近 completed-route oracle，却在 active chunk 内仍只能渐进补全并依赖 chunk boundary。二者不能只凭训练 PPL 选取。

\paragraph{TBPTT 截断。}
$U=1$ 的 512-token gradient horizon 保留 forward cache，却阻断后续 chunk 对早期 writer 的 credit assignment。增大 $U$ 会快速增加 activation memory，质量收益尚未测定。

\paragraph{压缩与 reader takeover。}
$r_K=r_V=32$ 相对 $d=768$ 是强压缩。block-specific canonicalizers 必须学习可组合坐标系；compiler position term、depth term、self-prefix 和 reader decoder 均可能形成 memory shortcut。需要通过 entropy、writer mass、self-memory concentration 与 public benchmark 联合诊断。

\paragraph{RoPE 限制低维 Key。}
Value read 已能在 latent space 精确聚合，但 position-dependent RoPE 阻止当前实现把 Key scoring 完全降到 $r_K$ 维。任何 latent-relative-RoPE 优化都必须先给出数学等价式和显式测试。

\paragraph{路由范围。}
v1 只实现 hard top-1，尚未支持 top-2 reader fusion；八个 free blocks、16 个最大 route steps 和 current routing guardrails 也尚未证明是规模最优解。

\section{实验协议}
当前 baseline 与 \DP{} C512/U1 5B runs 是 matched anchors，但不是 visibility ablation。后续机制验证按以下顺序进行：
\begin{enumerate}
  \item 在相同 $C,U$、global batch、数据与 token budget 下比较 \DP{} 与 \FB{}；
  \item 只在胜出的 visibility 上比较 $U=1,2,4$；
  \item 以 C128/U4 对 C512/U1，在相同 512-token gradient horizon 下测 prefill granularity；
  \item 从胜出配置分别比较 \code{bdre\_prefix} 与 \code{current\_step}、depth score on/off 和 compiler position on/off；
  \item hard top-$k$ 与温度只在 compiler entropy 或 writer domination 出现异常时启动。
\end{enumerate}
首轮每个 arm 使用 250M tokens，并在约 $1/3,2/3,1$ 进度执行 legacy validation、reasoning S600 与 public S600。只有综合能力、route/cache 稳定性和性能均成立的配置才晋级 2B；不构造每 arm 5B 的笛卡尔积。

\section{结论}
当前实现将自由路由中的 KV memory 分成三个清晰层次：block/head-specific attention 负责写入与读取的局部坐标系，canonical K/V latent 提供跨路径接口，\BDRE{} 依据 reader geometry 和可选 depth 规则把完整 writer route 编译成每 token 的 reader-specific object。\Serial{} 给出精确因果 oracle；\CPBC{} 通过 chunk 内 token 并行和 route-depth 串行实现可训练的近似 prefill，并用 DP/FB 明确暴露其历史可见性选择。stateful TBPTT 与手工 DDP 同步进一步保证 cache 生命周期和 optimizer update 一致。

系统的功能、因果性和分布式训练接口已经实现并通过测试；fused reader 已显著降低 activation memory 并提高吞吐，但系统级执行效率仍是 limitation。下一阶段的关键不是继续堆叠 routing loss，而是在不改变 \RCKV{}/\CPBC{} 语义的前提下减少 host synchronization、packing 与小 kernel 开销，再用受控 ablation 判断该记忆结构的真实能力收益。

\begin{thebibliography}{9}
\small
\bibitem{vaswani2017attention}
A. Vaswani et al., ``Attention Is All You Need,'' \emph{NeurIPS}, 2017.

\bibitem{su2021roformer}
J. Su et al., ``RoFormer: Enhanced Transformer with Rotary Position Embedding,'' \emph{Neurocomputing}, 2024; arXiv:2104.09864.

\bibitem{williams1990tbptt}
R. J. Williams and J. Peng, ``An Efficient Gradient-Based Algorithm for On-Line Training of Recurrent Network Trajectories,'' \emph{Neural Computation}, 1990.

\bibitem{fedus2022switch}
W. Fedus, B. Zoph, and N. Shazeer, ``Switch Transformers: Scaling to Trillion Parameter Models with Simple and Efficient Sparsity,'' \emph{JMLR}, 2022.
\end{thebibliography}

\end{document}
